% ===========================================================================
%  proof.tex — accepted artifact for problem `infinitely-many-primes`
%
%  Target (verbatim user request): "Prove that there are infinitely many primes."
%
%  Team:    math-team (natural-language workflow)
%  Published by: team-lead, mechanical publication of an accepted artifact.
%
%  SOURCE OF TRUTH — the body of this file is a VERBATIM copy, byte-for-byte,
%  of the accepted proof. No mathematical content was altered, reformatted,
%  or re-typeset in producing this file.
%
%      source:  nl/generator/proof-v1.md
%      SHA-256: 8a9cfda30bb3667442155db72547857cf6eac98f19687bf03bc236feb1260053
%
%  ACCEPTANCE BASIS
%      Independent, fresh, one-shot `math-nl-verifier` (CERTIFICATION mode);
%      the reviewer recomputed every artifact hash, matched the snapshot, and
%      returned VERDICT: PASS.
%      Review packet:  nl/reviews/proof-review-v1.md
%      SHA-256:        430b98d4365a885506c45e29f0e01eff584e58982342e9a73ce59bf5d54f7958
%
%  STATUS AND LIMITS — read this before citing the proof.
%      * This is an LLM-REVIEWED NATURAL-LANGUAGE PROOF. It is NOT a Lean
%        kernel/compiler-verified theorem. No formal (machine-checked)
%        verification of this statement has been performed.
%      * No TeX toolchain (pandoc / pdflatex / xelatex / lualatex / tectonic)
%        was installed at publication time, so this file was NOT compiled to
%        PDF and compilation was not tested. Compiling it would not confer
%        mathematical verification in any case; it is a typesetting step only.
%      * The body is included in a `verbatim` environment precisely so that
%        the accepted characters are preserved literally. A Unicode-capable
%        engine (lualatex or xelatex) would be required to compile it.
%
%      Provenance: the accepted artifact remains
%      nl/generator/proof-v1.md at the SHA-256 above; if this file and that
%      file ever disagree, the .md at that hash is authoritative.
% ===========================================================================
\documentclass[11pt]{article}
\usepackage[T1]{fontenc}
\usepackage{geometry}
\geometry{margin=1in}
\title{There are infinitely many primes}
\author{Team \texttt{math-team} (natural-language workflow)}
\date{Accepted 2026-09-24}
\begin{document}
\maketitle

\noindent\textbf{Note.} The material below is a verbatim copy of the accepted
natural-language proof \texttt{nl/generator/proof-v1.md}, SHA-256
\texttt{8a9cfda3\dots 260053}, independently reviewed in CERTIFICATION mode by a
fresh \texttt{math-nl-verifier} (packet \texttt{nl/reviews/proof-review-v1.md},
verdict \textbf{PASS}). It is an LLM-reviewed proof, \emph{not} a
Lean-kernel-verified theorem. It was not compiled here (no TeX toolchain), and
typesetting success would not constitute mathematical verification. See the
provenance block at the top of this file.

\begin{verbatim}
# Proof attempt — infinitely many primes

**Mode:** CERTIFICATION candidate (this artifact is a *proof attempt*; it carries no
proof weight until a fresh `math-nl-verifier` accepts it). The author does **not**
certify this result.

**Problem:** `infinitely-many-primes`. **Target (verbatim from `request.md`):**

> Prove that there are infinitely many primes.

**Inputs used:** `request.md`, `nl/sketcher/target-contract.md`,
`nl/sketcher/lemma-plan.md`, `nl/sketcher/obligations.md`.
Definitions, domain, and normalizations below are taken from the contract;
the lemmas `L1`–`L6` are the plan's statements but every one of them is proved
here from scratch. Nothing is cited as established merely because the plan
asserts it.

---

## 0. Target, domain, definitions, framework

### 0.1 The assertion and its two readings

Let
```
P := { p ∈ ℕ : p is prime }.
```
**Primary reading (cardinality).** `P` is infinite, i.e. `P` is not finite
(definition of finite in §0.4, Definition F).

**Working form (unboundedness).** `∀ n ∈ ℕ, ∃ p ∈ P with p > n`.

Per `target-contract.md` §1 these two readings are equivalent for subsets of ℕ;
that equivalence is *proved* here (Lemma L5 and Lemma L5c), not assumed. The
quantifier order of the working form is `∀n ∃p`; the swapped form
`∃p ∀n, p > n` is **false** and is not used anywhere. The argument below proves
the working form and then converts it to the primary reading through L5.

### 0.2 Domain (contract `N1`, `N2`)

`ℕ = {1, 2, 3, …}`, the positive integers. Primes are drawn from ℕ.

*Remark on `N2` (whether 0 ∈ ℕ).* The argument is written for `ℕ = {1,2,…}`.
If instead `0 ∈ ℕ`, then: (i) primes are still `> 1` by `D2`, so `0` is never
prime and the set `P` is literally unchanged; (ii) the working form at `n = 0`
follows from its value at `n = 1`, since any prime `p > 1` satisfies `p > 0`
(existence of a prime is Lemma L4); (iii) all divisibility occurrences below
have `d ≥ 1` or `n ≥ 1` (checked at each step). Hence both readings of ℕ yield
the same true statement, and the proof for the reading `ℕ = {1,2,…}` transfers.

### 0.3 Definitions (contract `D1`, `D2`)

**`D1` (divisibility).** For `a, b ∈ ℕ`: `a | b`  iff  `∃ c ∈ ℕ, b = a·c`.

**`D2` (prime; the normalized definition).** `p ∈ ℕ` is **prime** iff
```
p > 1   and   ∀ d ∈ ℕ, ( d | p  ⟹  d = 1 ∨ d = p ).
```
Consequences used: a prime satisfies `p > 1`, so `1` is not prime; and every
prime is `≥ 2`.

### 0.4 Finiteness (for the primary reading)

**Definition F.** A set `S` is **finite** iff either `S = ∅`, or there exist
`k ∈ ℕ` and a bijection `f : {1, …, k} → S`. A set is **infinite** iff it is not
finite. This matches `target-contract.md` §1 ("no `k ∈ ℕ` and no bijection
`{1,…,k} → P`"): since `P ≠ ∅` (e.g. `2 ∈ P`, proved in §1, E11), the empty
alternative never applies to `P`.

### 0.5 Ambient framework for ℕ

We use the standard model `ℕ = {1,2,3,…}`. The following are its standard
properties (each a theorem of the Peano axioms / the ordered semiring ℕ; **none
mentions primes**, so using them is not circular with respect to the target):

* **(AR1)** ℕ is a commutative semiring with multiplicative identity `1`;
  `1 ≤ a` for every `a ∈ ℕ`.
* **(AR2)** `≤` is a total order on ℕ (reflexive, antisymmetric, transitive,
  and total), compatible with arithmetic: `a ≤ b ⟹ a+c ≤ b+c` and `ac ≤ bc`.
* **(AR3)** `a < b ⟺ ∃ c ∈ ℕ, b = a + c`; and (discreteness) `a < b ⟹ a+1 ≤ b`.
* **(AR4)** Addition and multiplication cancel: `a+c = b+c ⟹ a = b`, and
  `ac = bc ⟹ a = b`.
* **(AR5)** Well-ordering: every nonempty `S ⊆ ℕ` has a least element.

All invocations of these are flagged as `O-T1` (AR5) and `O-T2`/`O-AR*` in the
ledger. No other infinity/factorization principle is used anywhere.

---

## 1. Elementary facts on ℕ, order, divisibility, and factorial

Every fact is proved from §0.5. Notation: `a > b` means `b < a`; `a ≥ b` means
`b ≤ a`.

**E1 (trichotomy / totality).** For `a, b ∈ ℕ`, exactly one of `a < b`, `a = b`,
`a > b` holds; in particular `a ≤ b` or `b ≤ a`.
*Proof.* `≤` is a total order (AR2), so `a ≤ b` or `b ≤ a`. Define `<` by
`a < b ⟺ a ≤ b ∧ a ≠ b` (AR3 gives the equivalent `∃c, b = a+c`). If `a ≤ b`
and `b ≤ a`, antisymmetry gives `a = b`; otherwise exactly one of `a < b`,
`b < a` holds. ∎

**E2 (discreteness, restated).** `a < b ⟹ a+1 ≤ b` (AR3).

**E3 (product dominates factors).** For `a, b ∈ ℕ`: `ab ≥ a` and `ab ≥ b`.
*Proof.* `b ≥ 1` (AR1). From `b ≥ 1`, multiply by `a`: `ab ≥ a·1 = a`
(AR2 + AR1). Symmetrically with `a ≥ 1`. ∎

**E4 (no factor `1` unless both are `1`).** If `a, b ∈ ℕ` and `ab = 1`, then
`a = b = 1`.
*Proof.* Suppose `a ≠ 1`. Since `a ≥ 1` (AR1) and `a ≠ 1`, E2 gives `a ≥ 2`.
Then `ab ≥ a ≥ 2 > 1` (E3), contradicting `ab = 1`. Hence `a = 1`; symmetrically
`b = 1`. ∎

**E5 (cancellation).** `ac = bc ⟹ a = b` (AR4). Also used: `a+c = b+c ⟹ a=b`
(AR4).

**E6 (divisibility basics).** For `a, b, c ∈ ℕ`:
(a) `a | a`;  (b) `1 | b`;
(c) *(transitivity)* `a | b` and `b | c` ⟹ `a | c`.
*Proof.* (a) `a = a·1`. (b) `b = 1·b`. (c) `b = a x`, `c = b y` with `x,y ∈ ℕ`;
then `c = a(xy)` and `xy ∈ ℕ`. ∎

**E7 (`|` implies `≤`).** If `a | b` then `a ≤ b`. Consequently, if `a | b` and
`b | a` then `a = b`.
*Proof.* `b = a c` with `c ∈ ℕ`, `c ≥ 1`; so `b = ac ≥ a·1 = a` by E3/AR2.
For the consequence apply E7 twice and antisymmetry. ∎

**E8 (only `1` divides `1`).** If `d ∈ ℕ` and `d | 1` then `d = 1`.
*Proof.* `1 = d·c` with `c ∈ ℕ`; E4 gives `d = 1`. ∎

**E9 (divisibility is linear under subtraction).** Let `a, b, d ∈ ℕ` with
`d | a`, `d | b`, and `a > b`. Then `d | (a − b)`, where `a − b` is the unique
`c ∈ ℕ` with `a = b + c` (it exists because `a > b`, and is unique by E5).
*Proof.* Write `a = d x`, `b = d y` with `x, y ∈ ℕ`. If `x ≤ y` then
`a = dx ≤ dy = b` (AR2), contradicting `a > b`. Hence by E1 `x > y`, so
`x = y + c'` for some `c' ∈ ℕ`. Then
`a = dx = d(y + c') = dy + dc' = b + dc'`.
By uniqueness of the difference, `a − b = dc'`, so `d | (a−b)`. ∎

**E10 (factorial).** Define `n!` for `n ∈ ℕ` by recursion:
`1! := 1`, and `(n+1)! := n! · (n+1)` for `n ≥ 1`. Then for all `n ∈ ℕ`:
(a) `n! ≥ 1`;  (b) `n! + 1 ≥ 2 > 1`.
*Proof.* (a) Induction: `1! = 1 ≥ 1`; and `(n+1)! = n!·(n+1) ≥ n! ≥ 1` by E3
and the induction hypothesis. (b) From (a), `n! + 1 ≥ 1 + 1 = 2 > 1`. ∎

**E11 (`2` is prime).** `2 ∈ P`; hence `P ≠ ∅`.
*Proof.* `2 > 1`. Let `d ∈ ℕ` with `d | 2`. By E7, `d ≤ 2`. Since `d ≥ 1` and
`1 ≤ d ≤ 2`, E2 (applied to `1 < d` when `d > 1`) gives `d = 1` or `d = 2`.
Hence every divisor of `2` in ℕ is `1` or `2`, so `2` is prime by `D2`. ∎

---

## 2. Lemma L1 — Every integer `> 1` has a prime divisor

**Statement (L1).** For every `m ∈ ℕ` with `m ≥ 2` there exists `p ∈ ℕ` such
that `p | m` and `p` is prime (`D2`).

**Proof.** Fix `m ∈ ℕ` with `m ≥ 2`. Define
```
S := { d ∈ ℕ : d | m  and  d > 1 }.
```
*Nonemptiness.* `m | m` by E6(a), and `m > 1` since `m ≥ 2`. Hence `m ∈ S`, so
`S ≠ ∅`. (Precondition of AR5 met: `S` is a nonempty subset of ℕ.)

*Choice of minimal element.* Apply **AR5 (well-ordering of ℕ; obligation
`O-T1`)**: `S` has a least element. Set
```
p := min S,
```
so `p ∈ S`; in particular `p | m` and `p > 1`.

*`p` is prime.* `p > 1` holds. Let `d ∈ ℕ` with `d | p`; we must show
`d = 1 ∨ d = p`. By E7, `d ≤ p`. By E1, either `d ≤ 1` — and then `d = 1`
since `d ≥ 1` — or `d > 1`. In the second case `d | p` and `p | m` give
`d | m` by **E6(c) (transitivity of `|`)**; together with `d > 1` this says
`d ∈ S`. As `p` is the least element of `S`, `p ≤ d`. With `d ≤ p` and
antisymmetry, `d = p`. Either way `d = 1` or `d = p`, so the `D2` divisor
condition holds and `p` is prime. ∎

**Non-circularity check (`O-L1`, keystone).** The proof uses **only**: the
definition `D1` of `|`; reflexivity (E6a) and transitivity (E6c) of `|`;
`| `-implies-`≤` (E7); E1 trichotomy; and **AR5 well-ordering**. It does **not**
use the infinitude of the primes (the target), does **not** use the Fundamental
Theorem of Arithmetic (existence *or* uniqueness), does **not** use Euclid's
lemma `p | ab ⟹ p | a ∨ p | b`, and does **not** assume `P` finite or infinite.
The minimal element `p` is produced directly by well-ordering, not by iterative
"extract a prime factor from a proper divisor" (which would presuppose L1). ∎

---

## 3. Lemma L2 — A prime `≤ n` divides `n!`

**Statement (L2).** For all `n ∈ ℕ` and every prime `p`: if `p ≤ n` then
`p | n!`.

**Proof.** Induction on `n ≥ 1` using the recursion of E10.

*Base `n = 1`.* If `p` is prime then `p > 1` (`D2`), so `p ≤ 1` is impossible
(by E1, `p ≤ 1` with `p ≥ 1` forces `p = 1`, contradicting `p > 1`). Hence the
statement is vacuously true for `n = 1`.

*Inductive step.* Assume L2 holds for `n ≥ 1`; prove it for `n+1`. Let `p` be
prime with `p ≤ n+1`. By **E1 (trichotomy)** either `p ≤ n` or `p > n`.
* If `p ≤ n`: the induction hypothesis gives `p | n!`. Also `n! | (n+1)!`,
  because the recursion gives `(n+1)! = n!·(n+1)`. By **E6(c) (transitivity of
  `|`)**, `p | (n+1)!`.
* If `p > n`: then `n < p ≤ n+1`. By **E2 (discreteness)**, `n < p` gives
  `n+1 ≤ p`; with `p ≤ n+1` and antisymmetry, `p = n+1`. Then
  `(n+1)! = n!·(n+1) = n!·p`, so `p | (n+1)!`.

The two cases are exhaustive (E1) and both yield `p | (n+1)!`. ∎

**Preconditions checked.** Factorial is defined for all `n ∈ ℕ` (E10); `p ∈ ℕ`
by `D2`. No hypothesis beyond `D1`, `D2`, and the factorial recursion is used.

---

## 4. Lemma L3 — `n! + 1` is coprime to `n!`

**Statement (L3).** For every `n ∈ ℕ`:
(i) `n! + 1 > 1`; and
(ii) every `d ∈ ℕ` with `d | n!` and `d | (n! + 1)` satisfies `d = 1`.

**Proof.** Fix `n ∈ ℕ`.

(i) By **E10(a)**, `n! ≥ 1`, so `n! + 1 ≥ 2 > 1`. ∎

(ii) Let `d ∈ ℕ` with `d | n!` and `d | (n! + 1)`. Put `a := n! + 1` and
`b := n!`. Then `a > b` (indeed `a = b + 1`), and `b, a, d ∈ ℕ`. Apply
**E9 (linearity of `|` under subtraction)** with these `a, b, d`: the
preconditions `d | a`, `d | b`, `a > b` all hold, so `d | (a − b)`. The
difference is `a − b = (n! + 1) − n! = 1` (the unique `c` with
`n! + 1 = n! + c` is `c = 1`, by E5). Hence `d | 1`, and **E8** gives `d = 1`. ∎

**Note.** No gcd or coprimality primitive is introduced; the statement is phrased
with `|` alone, as in the plan.

---

## 5. Lemma L4 — Unboundedness of the primes

**Statement (L4).** For every `n ∈ ℕ` there exists a prime `p ∈ ℕ` with `p > n`.

**Proof.** Fix `n ∈ ℕ`. Put
```
M := n! + 1.
```
By **E10(b)**, `M ≥ 2 > 1`, so the hypothesis `m ≥ 2` of **L1** holds for
`m := M`. Apply **L1**: there exists a prime `p ∈ ℕ` with `p | M`, i.e.
`p | (n! + 1)`.

Now apply **E1 (trichotomy)** to the pair `(p, n)`: either `p ≤ n` or `p > n`.
These two cases are exhaustive.
* **Case `p > n`.** Done: `p` is a prime with `p > n`.
* **Case `p ≤ n`.** Then `p` is prime and `p ≤ n`, so **L2** gives `p | n!`.
  Also `p | (n! + 1)`. Apply **L3(ii)** with `d := p` (legitimate since `p ∈ ℕ`):
  `p = 1`. But `p` is prime, so `D2` gives `p > 1`, contradicting `p = 1`. This
  case is therefore impossible.

Since the second case is contradictory, the first holds, and there is a prime
`p > n`. As `n ∈ ℕ` was arbitrary, `∀ n ∈ ℕ ∃ p prime, p > n`. ∎

**Hypotheses audit.** `L4` is a closed `∀n ∃p` statement: no finiteness
assumption on `P`, no bound on `p`, and no extra hypothesis is introduced. All
preconditions of L1, L2, L3 are checked above; the only order principle used is
E1 (from AR2/AR3).

*Remark (reading `N2`, `0 ∈ ℕ`).* If ℕ were taken to contain `0`, the case
`n = 0` is immediate from the case `n = 1` (there is a prime `p > 1 > 0`),
consistent with `target-contract.md` §3. The recursion-based proof above is
stated for `n ≥ 1`, which is exactly `ℕ = {1,2,…}`.

---

## 6. Lemma L5 — Finite subsets of ℕ are bounded above; and L5c (converse)

**Statement (L5).** For every `S ⊆ ℕ`: if `S` is finite (Definition F) then `S`
is bounded above, i.e. `∃ N ∈ ℕ, ∀ s ∈ S, s ≤ N`.

**Proof.** Two cases.

*Case `S = ∅`.* Take `N := 1 ∈ ℕ`. The condition `∀ s ∈ ∅, s ≤ 1` is vacuous. ∎

*Case `S ≠ ∅` finite.* By Definition F there are `k ∈ ℕ` and a bijection
`f : {1,…,k} → S`. We prove the following by induction on `k ≥ 1`:

> **(L5-g)** For every `k ∈ ℕ`, every `S ⊆ ℕ` admitting a bijection
> `f : {1,…,k} → S` has a **greatest element**: `∃ g ∈ S, ∀ s ∈ S, s ≤ g`.

*Base `k = 1`.* `S = {f(1)}`; take `g := f(1) ∈ S ⊆ ℕ`. For `s ∈ S`, `s = f(1)`
so `s ≤ g`. ✓

*Inductive step (`k → k+1`).* Assume L5-g for `k ≥ 1`. Let `S ⊆ ℕ` and
`f : {1,…,k+1} → S` be a bijection. Set `S' := S \ { f(k+1) }`; the restriction
`f|_{\{1,…,k\}}` is a bijection `{1,…,k} → S'`. By the induction hypothesis `S'`
has a greatest element `g'`. Define `g := max(g', f(k+1))`, the larger of two
elements of ℕ (so `g ∈ ℕ`). For any `s ∈ S`: either `s = f(k+1)`, giving
`s = f(k+1) ≤ g`; or `s ∈ S'`, giving `s ≤ g' ≤ g`. Hence `g` is a greatest
element of `S`. ✓

Applying L5-g to the given `S` and `f` yields a greatest element `g ∈ S ⊆ ℕ`.
Take `N := g ∈ ℕ`. For every `s ∈ S`, `s ≤ g = N`. Hence `S` is bounded above. ∎

**Contrapositive form (used for the target).**

> **(L5\*)** For every `S ⊆ ℕ`: if `S` is *unbounded above* — i.e.
> `∀ n ∈ ℕ ∃ s ∈ S, s > n` — then `S` is infinite (not finite).

*Proof.* `L5*` is the literal contrapositive of L5: "finite ⟹ bounded above"
is logically equivalent to "not bounded above ⟹ not finite"; "not bounded
above" is the displayed unboundedness condition, and "not finite" is the
definition of infinite. ∎

**Statement (L5c, converse; `O-L5c`, non-load-bearing).** If `S ⊆ ℕ` is bounded
above then `S` is finite.

*Proof.* Suppose `∃ N ∈ ℕ, ∀ s ∈ S, s ≤ N`. Since every `s ∈ S ⊆ ℕ` satisfies
`s ≥ 1`, we get `S ⊆ {1,…,N}`. We show by induction on `N ≥ 1` that **every**
subset of `{1,…,N}` is finite.
* Base `N = 1`: a subset of `{1}` is `∅` or `{1}`; `∅` is finite by Definition F,
  and `{1}` is finite via the identity bijection `{1} → {1}`.
* Step `N → N+1`: let `T ⊆ {1,…,N+1}`. If `N+1 ∉ T` then `T ⊆ {1,…,N}`, finite
  by the induction hypothesis. If `N+1 ∈ T`, write `T = T' ∪ {N+1}` with
  `T' ⊆ {1,…,N}`, finite by the induction hypothesis; if `T' = ∅` then
  `T = {N+1}`, finite via `1 ↦ N+1`; otherwise `T'` has a bijection
  `f: {1,…,k} → T'` (some `k ∈ ℕ`) and extending it by `k+1 ↦ N+1` is a
  bijection `{1,…,k+1} → T`, so `T` is finite.
Hence `S ⊆ {1,…,N}` is finite. ∎

**Remark.** L5 and L5c together give, for every `S ⊆ ℕ`: `S` is infinite `⟺`
`S` is unbounded above. This is exactly the equivalence asserted in
`target-contract.md` §1 / ambiguity `A1`; it is now proved, so the reading
choice is discharged rather than assumed.

---

## 7. Lemma L6 — Assembly: the original target

Let `P := { p ∈ ℕ : p prime }` (`D2`).

1. **`P ⊆ ℕ`.** Immediate from the definition of `P`. (`O-B1`.)
2. **Working form.** By **L4**, for each `n ∈ ℕ` there is a prime `p ∈ ℕ` with
   `p > n`. That `p` satisfies `p ∈ P` and `p > n`. Hence
   `∀ n ∈ ℕ, ∃ p ∈ P with p > n`, i.e. `P` is unbounded above. The `∀n ∃p`
   order of L4 is preserved verbatim; the swapped form `∃p ∀n` is **not**
   produced or used. (`O-B3`.)
3. **Unbounded ⟹ infinite.** Apply **L5\*** with `S := P` (legitimate: `P ⊆ ℕ`
   by step 1, and `P` is unbounded above by step 2). Conclusion: `P` is infinite.
4. **Reading bridge.** By `target-contract.md` §1 the primary reading of
   "there are infinitely many primes" **is** the statement "`P` is infinite".
   Step 3 therefore establishes the original target. (`O-B2`.)

**Conclusion.** `P = { p ∈ ℕ : p prime }` is infinite; equivalently,
`∀ n ∈ ℕ ∃ p ∈ ℕ (p prime ∧ p > n)`. ∎

---

## 8. Global audit

**Hypothesis audit.** No hypothesis is added anywhere. The target is an
unconditional closed statement. In particular:
* no finiteness assumption on `P` is made in L4 (the contradiction-style Route B
  of the plan is deliberately *not* used);
* no genericity, nonzero, or regularity hypothesis is inserted: all uses of
  divisibility occur with arguments in ℕ, and the "`d ≥ 1`"/"`d > 1`" conditions
  are forced by `D2`, E1, and E10, as shown;
* the ambient principles (AR1)–(AR5) are standard properties of ℕ and do not
  mention primes.

**Circularity audit.** The dependency graph is
`D1,D2,AR* → {L1, L2, L3, L5(-g)} → L4 → L5* → L6`.
No lemma uses the target. Specifically:
* L1 uses well-ordering of ℕ, `D1`, and transitivity of `|` only (§2);
* L2 uses the factorial recursion and transitivity of `|` only (§3);
* L3 uses subtraction-linearity of `|` (E9) and E8 only (§4);
* L4 uses L1, L2, L3 and trichotomy only (§5);
* L5/L5-g use Definition F, induction, and `max` on ℕ only (§6);
* the assembly uses L4, L5\*, and the definitional reading bridge only (§7).
Explicitly **not** used: infinitude of primes, FTA (existence or uniqueness),
Euclid's lemma, Dirichlet/Bertrand/analytic input.

**Quantifier audit.** The proven statement is `∀ n ∈ ℕ ∃ p ∈ P, p > n`
(unboundedness) and the cardinality statement "`P` is infinite". The false
swapped form is never asserted.

**Boundary audit.** `1` is not prime (`D2`, `p > 1`); `2` is the least prime
(E11); the empty and finite cases of "infinitely many" are handled by L5/L5c and
the reduction to the non-finite reading.

---

## 9. Status and unresolved issues

Every step above is discharged by an explicit argument from `D1`, `D2`, and
(AR1)–(AR5); the obligation ledger `obligations.md` maps each named obligation
of the plan to its discharge. **No unresolved gap remains in this proof
attempt**, apart from the following caveats, which are limitations of scope
rather than gaps:

* The proof is written for the reading `ℕ = {1,2,…}` (`N1`); the transfer to a
  reading in which `0 ∈ ℕ` is argued in §0.2/§5 and rests only on `p > 1` for
  primes.
* `L5c` / `O-DEF2` (equivalence `D2 ⟺ D3`, and `bounded ⟺ finite`) are supplied
  for completeness / to discharge ambiguity `A1`; `O-DEF2` is not needed by the
  route and is not proved here (it is recorded as an optional completeness item
  in `obligations.md`).
* Acceptance is **not** claimed by the author. A fresh independent
  `math-nl-verifier` must check this candidate in a separate task.

\end{verbatim}

\end{document}
